% Esquisse d'un programme (cote 119, copy 2): archive cote 119, site page 81 (PDF page 82), typescript page 3.
% typescript. Draft compared: esquisse/pages/esq_119_005.tex (copy 1, site 5).
% AI-transcribed by Claude Opus 5.5 (claude-opus-5-5, adaptive thinking, effort high), Message Batches API, batch esq119_full (msgbatch_015xqDHkTumDysRF1qngZhRJ), 2026-11-13; not yet checked by eye.
\scanpage{119/81}
\mspage{3}

1970, mon attention systématiquement était dirigée vers les objets
de généralité maximale, afin de dégager un langage d'ensemble adéquat
pour le monde de la géométrie algébrique, et que je ne m'attardais
sur les courbes algébriques que dans la stricte mesure où cela s'avé-%
rait indispensable (notamment en cohomologie étale) pour développer
des techniques et énoncés ``passe-partout'' valables en toutes dimen-%
sions et en tous lieux (j'entends, sur tous schémas de base, voire
tous topos annelés de base ...), me voici donc ramené, par le tru-%
chement d'objets si simples qu'un enfant peut les connaître en jouant,
aux débuts et origines de la géométrie algébrique, familiers à Riemann
et à ses émules !

Depuis environ 1975, c'est donc la géométrie des surfaces
(réelles), et à partir de 1977 les liens entre les questions de géo-%
métrie des surfaces et la géométrie algébrique des courbes algébri-%
ques définies sur des corps tels que $\mathbb{C}$, $\mathbb{R}$ ou des extensions de type
fini de $\mathbb{Q}$, qui ont été ma principale source d'inspiration, ainsi
que mon fil conducteur constant. C'est avec surprise et avec émer-%
veillement qu'au fil des ans je découvrais (ou plutôt, sans doute,
redécouvrais) la richesse prodigieuse, réellement inépuisable, la
profondeur insoupçonnée de ce thème, d'apparence si anodine. Je crois
y sentir un point névralgique entre tous, un point de convergence
privilégié des principaux courants d'idées mathématiques, comme aussi
des principales structures et des visions des choses qu'elles expri-%
ment, depuis les plus spécifiques (tels les anneaux $\mathbb{Z}$, $\mathbb{Q}$, $\bar{\mathbb{Q}}$, $\mathbb{R}$, $\mathbb{C}$,
ou le groupe $S\hand{\ell}(2)$ sur l'un de ces anneaux, ou les groupes algébri-%
ques réductifs généraux) aux plus ``abstraits'', telles les ``multipli-%
cités'' algébriques, analytiques complexes ou analytiques réelles.
(Celles-ci s'introduisent naturellement quand il s'agit d'étudier
systématiquement des ``variétés de modules'' pour les objets géométri-%
ques envisagés, si on veut dépasser le point de vue notoirement in-%
suffisant des ``modules grossiers'', qui revient à ``tuer'' bien malen-%
contreusement les groupes d'automorphismes de ces objets .) Parmi ces
multiplicités modulaires, ce sont celles de Mumford-Deligne pour les
courbes algébriques ``stables'' de genre $g$, à $\nu$ points marqués, que

\margin{81}
